\section{Planar equivalent layers}
\label{sec:eqlayers}

The integral equations representing the generalized solutions of Dirichlet's
and Neumann's problems (equations \ref{eq:U-Dirichlet-cross} and \ref{eq:U-Neumann-cross})
can be alternatively interpreted in terms of their kernels.
Note that, according to equation \ref{eq:U-Dirichlet-cross}, the potential 
$U(\mathbf{r})$ is analogous to the vertical component of the gravitational 
attraction produced at all points $\mathbf{r}$ located above $z''_{0}$ by a fictitious 
planar layer of monopoles defined on $z''_{0}$, with mass distribution
$\tau(\mathbf{r}'')$.
This equation also represents a solution of the planar equivalent-layer integral
(equation \ref{eq:integral-equation-eqlayer}) for the particular case in which 
the potential field $d(\mathbf{r})$ is the gravitational potential 
$U(\mathbf{r})$ (equation \ref{eq:scalar-potential-U}) due to the
3D sources and the kernel $a(\mathbf{r}, \mathbf{r}'')$ is defined by 
$\partial_{z} \Psi(\mathbf{r}, \mathbf{r}'')$. In this case, 
$p(\mathbf{r}'') = \tau(\mathbf{r}'')$, which represents (i) the actual gravitational
potential produced by the 3D sources on a plane $z''_{0}$ located above the sources
or (ii) a fictitious gravitational potential obtained by downward continuation to points
inside the sources.

Equation \ref{eq:U-Neumann-cross}, instead, shows that  
$U(\mathbf{r})$ is analogous to the gravitational potential produced at all points
$\mathbf{r}$ located above $z''_{0}$ by a fictitious planar layer of monopoles defined on
$z''_{0}$, with mass distribution $\upsilon(\mathbf{r}'')$.
Besides, it defines another solution of the planar equivalent-layer integral
(equation \ref{eq:integral-equation-eqlayer}) for the case in which the
potential field $d(\mathbf{r})$ is the gravitational potential 
$U(\mathbf{r})$ (equation \ref{eq:scalar-potential-U}) due to the 3D sources
and the kernel is given by $a(\mathbf{r}, \mathbf{r}'') = \Psi(\mathbf{r}, \mathbf{r}'')$.
In this case, the continuous physical property distribution $p(\mathbf{r}'')$ is defined
by $\upsilon(\mathbf{r}'')$, which represents (i) the actual vertical component of the
gravitational attraction produced by the 3D sources on a 
plane $z''_{0}$ located above the sources or (ii) a fictitious gravity disturbance 
obtained by downward continuation to points inside the sources.


\subsection{Planar equivalent layer of monopoles applied to gravity data}
\label{subsec:eqlayer_monopoles}


The integral equations representing the generalized solutions of Dirichlet's
and Neumann's problems (equations \ref{eq:U-Dirichlet-cross} and \ref{eq:U-Neumann-cross})
represent solutions for the planar equivalent-layer integral for the case in which
the potential field $d(\mathbf{r})$ is the gravitational potential 
$U(\mathbf{r})$ (equation \ref{eq:scalar-potential-U}) due to the 3D sources.
In practical situation, however, the potential-field measurements are represented
by the partial derivatives of first and second orders of $U(\mathbf{r})$
(equations \ref{eq:1st-deriv-U} and \ref{eq:2nd-deriv-U}).
Here, we show the planar equivalent layers associated with these derivatives can be
obtained from the generalized solution of Neumann's problem 
(equation \ref{eq:U-Neumann-cross}).

Consider the partial derivatives of first and second orders of $U(\mathbf{r})$ obtained
by differentiating equation \ref{eq:U-Neumann-cross}:
\begin{equation}
\partial_{\alpha} U(\mathbf{r}) =  
\int\limits_{-\infty}^{\infty} \int\limits_{-\infty}^{\infty}
\upsilon(\mathbf{r}'') \:\: 
\partial_{\alpha} \Psi(\mathbf{r}, \mathbf{r}'') \: dx''dy'' \: , 
\quad z < z''_{0} \: ,
\quad \alpha = x, y, z \: ,
\label{eq:1st-deriv-U-Neumann}
\end{equation}
and
\begin{equation}
\partial_{\alpha\beta} U(\mathbf{r}) =  
\int\limits_{-\infty}^{\infty} \int\limits_{-\infty}^{\infty}
\upsilon(\mathbf{r}'') \:\: 
\partial_{\alpha\beta} \Psi(\mathbf{r}, \mathbf{r}'') \: dx''dy'' \: , 
\quad z < z''_{0} \: ,
\quad \alpha, \beta = x, y, z \: .
\label{eq:2nd-deriv-U-Neumann}
\end{equation}
Note that the derivatives are computed with respect to the coordinates
of $\mathbf{r}$ and, consequently, act only on the kernel $\Psi(\mathbf{r}, \mathbf{r}'')$
and not on the function $\upsilon(\mathbf{r}'')$ 
(equation \ref{eq:U-Neumann-cross}).

Note that equation \ref{eq:1st-deriv-U-Neumann} represents a solution of the
planar equivalent layer integral for the case in which the potential field 
$d(\mathbf{r})$ is the vertical component $\partial_{z} U(\mathbf{r})$ of the
gravitational attraction due to the 3D sources, which is commonly used to
approximate gravity disturbance data \cite[e.g.,][ p. 93]{hofmann-wellenhof-moritz2005}.
In this case, the kernel and physical property distribution are given by 
$a(\mathbf{r}, \mathbf{r}'') = \partial_{z} \Psi(\mathbf{r}, \mathbf{r}'')$ and
$p(\mathbf{r}'') = \upsilon(\mathbf{r}'')$, respectively.
This result shows that there exists a continuous planar
equivalent layer of monopoles that exactly reproduces the gravity disturbance
data above the sources.
\citet{dampney1969} was the pioneer in using this result to interpolate gravity data and
a considerable amount of literature has been published on this topic since then.
This solution $p(\mathbf{r}'')$ can also be directly deduced from the
\textit{upward-continuation integral}
\citep[e.g., ][]{peters1949, roy1962, bhattacharyya1967, henderson1970, gunn1975}, 
which represents the solution of Dirichlet’s problem on a plane
(equation \ref{eq:U-Dirichlet}), or from the generalized solution presented here
(equation \ref{eq:U-Dirichlet}-cross).

Consider now the case in which the potential-field data $d(\mathbf{r})$ in the 
planar equivalent layer integral (equation \ref{eq:integral-equation-eqlayer}) 
are gravitational gradient tensor components. These data are usually described as the
partial derivatives of second order of the gravitational potential. Hence, by using the
second derivatives defined by equation \ref{eq:2nd-deriv-U-Neumann}, we obtain for the 
$\alpha\beta$-component of gravity tensor that
$d(\mathbf{r}) = \partial_{\alpha\beta} U(\mathbf{r})$ and 
$a(\mathbf{r}, \mathbf{r}'') = \partial_{\alpha\beta} \Psi(\mathbf{r}, \mathbf{r}'')$, 
so that the continuous equivalent layer has a closed-form solution 
$p(\mathbf{r}'') = \upsilon(\mathbf{r}'')$.
This result shows that the same continuous planar equivalent layer of monopoles that 
reproduces the observed gravity disturbance data also reproduces the 
gradient tensor data above the sources.
This result has been used by the scientific community for processing gravity-gradient
data \citep{barnes-lumley_2011} or converting gravity disturbance into the
gradient tensor \citep{piauilino_etal2019_seg, li_etal2021}, for example.
However, researchers have not treated this topic in much detail and it is usually assumed
to be true without any proof.

\subsection{Planar equivalent layer of dipoles applied to magnetic data}
\label{subsec:eqlayer_dipoles}

We may also combine the partial derivatives of first order 
(equation \ref{eq:1st-deriv-U-Neumann}) to compute the
following directional derivative of first order:
\begin{equation}
\partial_{h} U(\mathbf{r}) = 
\int\limits_{-\infty}^{\infty} \int\limits_{-\infty}^{\infty}
\upsilon(\mathbf{r}'') \:\: 
\partial_{h} \Psi(\mathbf{r}, \mathbf{r}'') \:\: dx''dy'' \: , 
\quad z < z''_{0} \: ,
\quad \alpha, \beta = x, y, z \: ,
\label{eq:oblique-1st-deriv-U-Neumann}
\end{equation}
where 
\begin{equation}
\partial_{h} \Psi(\mathbf{r}, \mathbf{r}'') = 
\nabla \Psi(\mathbf{r}, \mathbf{r}'')^{\top} \hat{\mathbf{h}} 
\label{eq:oblique-1st-deriv-Psi}
\end{equation}
is the directional derivative of the inverse distance function 
$\Psi(\mathbf{r}, \mathbf{r}'')$ (equation \ref{eq:Psi}) along the
unit vector $\hat{\mathbf{h}}$ and
$\nabla \Psi(\mathbf{r}, \mathbf{r}'')$ is the gradient of $\Psi(\mathbf{r}, \mathbf{r}'')$
with respect to the coordinates of the observation point $\mathbf{r}$.

Finally, we may use partial derivatives of second order 
(equation \ref{eq:2nd-deriv-U-Neumann}) to obtain the following directional derivative
of second order:
\begin{equation}
\partial_{th} U(\mathbf{r}) = 
\int\limits_{-\infty}^{\infty} \int\limits_{-\infty}^{\infty}
\upsilon(\mathbf{r}'') \:\: 
\hat{\mathbf{t}}^{\top} \mathbf{H}_{\Psi}(\mathbf{r}, \mathbf{r}'') \hat{\mathbf{h}} 
\:\: dx''dy'' \: , \quad z < z''_{0} \: ,
\quad \alpha, \beta = x, y, z \: ,
\label{eq:oblique-2nd-deriv-Theta-Neumann}
\end{equation}
where
\begin{equation}
\mathbf{H}_{\Psi}(\mathbf{r}, \mathbf{r}'') = 
\begin{bmatrix}
\partial_{xx} \Psi(\mathbf{r}, \mathbf{r}'') & 
\partial_{xy} \Psi(\mathbf{r}, \mathbf{r}'') & 
\partial_{xz} \Psi(\mathbf{r}, \mathbf{r}'') \\
\partial_{xy} \Psi(\mathbf{r}, \mathbf{r}'') & 
\partial_{yy} \Psi(\mathbf{r}, \mathbf{r}'') & 
\partial_{yz} \Psi(\mathbf{r}, \mathbf{r}'') \\
\partial_{xz} \Psi(\mathbf{r}, \mathbf{r}'') & 
\partial_{yz} \Psi(\mathbf{r}, \mathbf{r}'') & 
\partial_{zz} \Psi(\mathbf{r}, \mathbf{r}'')
\end{bmatrix} \: ,
\label{eq:H-Psi}
\end{equation}
is the Hessian matrix of the inverse distance function 
$\Psi(\mathbf{r}, \mathbf{r}'')$ (equation \ref{eq:Psi}) defined in terms of
the second derivatives $\partial_{\alpha\beta} \Psi(\mathbf{r}, \mathbf{r}'') \equiv 
\frac{\partial^{2}}{\partial\alpha \, \partial\beta} \Psi(\mathbf{r}, \mathbf{r}'')$,
$\alpha = x, y, z$, $\beta = x, y, z$, with respect to the coordinates of
the observation point $\mathbf{r}$.